\chapter{高频、微观结构与执行：从订单流到事件时间研究}

\begin{itemize}
  \item \textbf{直接先修：}上册第 7、8、9、11、12、14、15、16、17 章；路线诊断见第 \ref{route:microstructure} 节。
  \item \textbf{学习产出：}定义事件时间与点过程，手工回放 FIFO 订单簿，以动态规划构造库存感知执行，并交付可复现事件研究包。
  \item \textbf{研究边界：}合成到达和简化订单簿不含真实交易所协议、延迟、隐藏流动性或市场冲击反馈。
\end{itemize}

\section{事件时间、强度与似然}

日历时间每秒推进一次，事件时间只在报价、委托、撤单或成交发生时推进。令 $N(t)$（离散记号也写作 $N_t$）为截至 $t$ 的事件数，条件强度定义为
\[
 \lambda(t\mid\mathcal F_{t^-})
 =\lim_{h\downarrow0}\frac{\mathbb P(N(t+h)-N(t)=1\mid\mathcal F_{t^-})}{h}.
\]
它是给定事件发生前信息的局部到达率，不是某个固定窗口计数的同义词。

齐次 Poisson 过程在区间 $[0,T]$ 观察到 $n$ 次事件的似然与到达时刻无关：
\[
 \ell(\lambda)=n\log\lambda-\lambda T,
 \qquad \widehat\lambda=\frac nT.
\]
本章四个到达间隔为 $(0.5,1,1.5,1)$，所以 $T=4$、$n=4$，强度为 1，对数似然为 $-4$。这两个数来自手算，不由模拟器生成。

市场事件常聚集。指数核 Hawkes 过程可写成
\[
 \lambda_t=\mu+\sum_{t_i<t}\alpha e^{-\beta(t-t_i)}.
\]
分枝比 $\alpha/\beta<1$ 是平稳版本的关键条件；若不小于 1，长期事件率可能爆炸。本章显式拒绝不稳定参数，但不以一次看似平稳的有限样本替代理论条件。

齐次强度还可能被日内季节性破坏。若两个等长窗口计数为 8 与 2，最大最小比超过声明阈值，程序以“日内季节性”拒绝常强度模型。更合理的流程是先估计确定性季节曲线，再对去季节后的残差检验聚集性，而不是把开收盘高峰错误归为自激。

\begin{diagnostic}
事件相关不等于收益可预测。点过程拟合必须分别检查时间变换残差、稳定性、样本外似然和经济用途；训练内似然提高不能单独证明交易价值。
\end{diagnostic}

\section{订单簿状态、事件与价格时间优先}

简化订单簿由买卖两侧的限价单队列组成。每条订单保存唯一编号、方向、价格、剩余数量与到达序号。必须维持四个不变量：数量为正；订单编号唯一；买单按高价优先、卖单按低价优先；同价位按到达序号 FIFO。

事件协议只包含三类动作：
\begin{enumerate}[leftmargin=*]
  \item \textbf{add：}在指定价位队尾加入限价单；
  \item \textbf{cancel：}从指定活动订单减少不超过剩余量的数量；
  \item \textbf{market：}从对手方最优价开始，按价格时间优先逐笔成交，允许部分成交。
\end{enumerate}

手工回放先在买价 100 加入 \verb|b1:5| 和 \verb|b2:4|，在卖价 101 加入 \verb|a1:6|。卖出市价单 7 先吃完 \verb|b1|，再成交 \verb|b2| 的 2；随后撤销 \verb|b2| 的 1，故买一剩 1。买入市价单 4 部分成交 \verb|a1|，卖一剩 2。最终买一剩 1、卖一剩 2，累计成交 11；成交编号顺序是 \verb|b1,b2,a1|。

错误撤单不能静默截断。若撤单量超过活动剩余量，状态与上游事件至少有一个错误，程序必须拒绝并保留事件上下文。同理，市场单超过显示流动性时，本章选择拒绝；真实交易所可能继续扫更深档位，协议必须按市场制度明确实现。

Python 字典很容易表达价位，却不会自动提供稳定 FIFO、唯一订单编号和原子状态更新。数据结构选择必须服务撮合不变量；测试也应观察成交顺序和剩余队列，而非直接断言私有容器布局。

\section{最优执行、状态控制与边界条件}

设尚待执行库存为 $x_t$，本期交易量为 $q_t$，状态转移为 $x_{t+1}=x_t-q_t$。本章用离散整数模型
\[
 \min_{q_0,\ldots,q_{T-1}}
 \sum_{t=0}^{T-1}\left[\eta q_t^2+\gamma x_{t+1}^2\right],
 \quad x_0=X,\quad x_T=0,
\]
其中 $\eta q_t^2$ 是临时冲击，$\gamma x_{t+1}^2$ 惩罚剩余库存风险。每期还受 $0\le q_t\le q_{\max}$ 和显示流动性限制。

动态规划递推为
\[
 V_t(x)=\min_q\left\{\eta q^2+\gamma(x-q)^2+V_{t+1}(x-q)\right\},
\]
边界 $V_T(0)=0$，未清库存的终态成本为 $+\infty$。对 $X=6$、$T=3$、$\eta=1$，若 $\gamma=0$，凸性给出等量 TWAP $(2,2,2)$，冲击成本为 12。取 $\gamma=0.3$、每期最多 3，则枚举与动态规划都得到 $(3,2,1)$，总目标为 17。

前置交易不是免费的“更聪明”。它以更大即时冲击换取更小持仓风险。报告应分解冲击、价差、机会成本和库存项，而不是只列总目标。参数也应来自可估计或可审计的协议；用回测收益倒推到最好看的风险厌恶系数属于研究自由度。

零冲击会让大单执行看似免费，无限显示流动性会删除参与率约束，未来价格可见则直接构成时间泄漏。本章把三者作为独立失败类别。真实执行还需处理永久冲击、随机成交、订单选择、延迟、部分完成和收盘边界。

\begin{note}[库存控制]
做市问题通常把买卖报价作为控制，以价差收益对抗库存风险与逆向选择。本章的单边清仓模型只提供状态、控制、成本和终端边界的最小桥接，不应被描述为完整做市器。
\end{note}

\section{可复现仿真与统计评估}

仿真器必须固定随机数生成器、种子、事件协议和停止规则。本章用强度 2 的指数到达间隔生成 2000 个事件；理论均值和方差都为 $0.5$ 与 $0.25$。固定种子样本的均值约为 $0.494991$、方差约为 $0.237309$，均值 95\% 区间覆盖理论值 0.5。

一次覆盖不验证置信区间的长期覆盖率。更完整的实验应重复多个种子，报告估计偏差、方差、区间覆盖与计算预算。若研究策略收益，还要把同一事件流送入简单基线和候选策略，使用配对差异减少随机源噪声。

事件回放与随机仿真承担不同职责：前者用手工台账验证状态机，后者检验统计性质。不能用“模拟均值接近理论值”替代 FIFO、撤单和部分成交测试，也不能用一条正确回放证明强度模型合适。

延迟必须至少分为行情、决策、发单、网络与撮合延迟；冲击需区分临时与永久部分；库存报告要同时给路径峰值和终值。交易所的价格时间优先、最小价位、涨跌停、自成交保护和撤单规则不同，研究包必须声明适用制度。

\section{Route Capstone：事件时间执行研究包}

Capstone 从冻结 fixture 运行点过程手算、种子仿真、订单簿回放、TWAP 基线与动态规划执行。六类失败分别守住季节性、Hawkes 稳定性、撤单不变量、正冲击、有限流动性和信息时点。

冻结摘要见 \path{reports/lower-ch06-summary.md}。陌生审计者从干净目录运行后，应得到相同成交顺序、剩余队列、执行序列与统计报告；这仍不代表低延迟工程、生产做市或实盘成交能力。

\begin{implementationnote}
\begin{lstlisting}[language=bash]
uv run python notebooks/lower/ch06_microstructure.py \
  evidence/lower-ch06/oracle.json
\end{lstlisting}
\end{implementationnote}

\section{四级问题}
\begin{enumerate}[leftmargin=*]
  \item \textbf{口述概念：}事件时间与日历时间有什么差别？条件强度依赖哪份信息？
  \item \textbf{口述概念：}Hawkes 分枝比为何必须小于 1？日内季节性为何会伪装成聚集？
  \item \textbf{笔试推导：}推导齐次 Poisson 强度的极大似然估计。
  \item \textbf{笔试推导：}写出库存感知执行的 Bellman 递推和终端边界。
  \item \textbf{数值编程：}扩展手工回放，加入同价第三张订单与部分撤单，核对 FIFO 成交序列。
  \item \textbf{数值编程：}枚举六股三期执行方案，验证 TWAP 与库存惩罚解。
  \item \textbf{研究判断：}如何区分到达模型错设、订单簿状态机错误与策略表现不佳？
  \item \textbf{研究判断：}为真实市场补充延迟、冲击、制度差异和不可声称能力的报告协议。
\end{enumerate}

